\section{Scope, provenance, and conventions}\label{sec:scope}

\subsection{What this continuation establishes}
The requested direction is exact special-function identities, rather than
inequalities or numerical nonreduction.  The present development connects
five objects: gamma ratios, generalized harmonic numbers, spectral
Hurwitz-zeta derivatives, generalized Stieltjes constants, and polylogarithmic
endpoint counterterms.  Asymptotic expansions enter only to prove that
certain explicitly displayed differences are genuinely summable and
that differentiation is legitimate.  The final statements are identities.

There are three main mechanisms.  First, the poles of Gauss's summation
formula cancel the poles of the matching Hurwitz-zeta counterterms.
Second, the asymptotic coefficients satisfy a parameter differential law,
which collapses every zeroth-order resonance to a short digamma expression.
Third, Taylor coefficients preserve the cancellation and produce Stieltjes
constants at precisely the colliding pole.  A different, balanced parameter
variation generates the all-weight harmonic-tail identity of
Theorem~\ref{thm:tail}.

The initial canonical repository README reports an exact proof of $S_4$
and keeps $S_6$ and the revised $S_8$ conjectural.  It also distinguishes
analytic proofs, finite exact certificates, and small numerical residuals
\cite{repo-readme}.  Those distinctions are retained here.  We do not
reprove $S_4$, promote either remaining Gaussian conjecture, or infer any
independence statement from numerical searches.

The three inspected incoming packages on shifted Hurwitz jets and
Stieltjes correlations already develop circular convolution, contact terms,
and primitive ladders.  Our construction instead starts from a
hypergeometric coefficient sequence.  The full incoming directory was
changing during the investigation; the audit is not exhaustive.  The latest
observed repository head was
\src{e0d9463bdee9685dfb1dddb819059cc738540c57}.
Individual inspected files and limitations are recorded in
\src{notes/SOURCE_AUDIT.md}.  A small wording correction was rechecked at
that exact commit; no repository file was modified.

\subsection{Classical inputs and the priority boundary}
Gauss's ${}_2F_1(1)$ evaluation, Stirling--Bernoulli gamma-ratio expansions,
and the defining identities for Hurwitz zeta are classical
\cite{dlmf-gauss,dlmf-gamma,dlmf-hurwitz}.
B\"uhring's treatment of hypergeometric partial sums includes zero and
negative integer excess and their constant terms \cite{buhring2004}.
In particular, the zeroth-order evaluation below belongs to that classical
constant-term setting; its short coefficient proof is an independent
organization, not a claim to have discovered hypergeometric finite parts.
Our coefficients $A_j$ are \emph{not} B\"uhring's coefficients
$A_k^{(p)}$: the latter encode a different continuation expansion.
The present fixed-shift subtraction basis, all-order Stieltjes coefficient
formula, degeneracy-safe implementation, and combined primitive calculus
are the specific development supplied to the repository.

Parameter differentiation of classical summations is likewise an
established source of harmonic identities.  We give complete proofs of
our particular all-weight formulas, but do not claim that every
specialization is absent from the literature.  A full priority audit of
individual binomial sums would require comparison with a much larger
hypergeometric literature.  No famous open problem is declared solved.

\subsection{Notation and analytic domains}
Unless a local complex extension is explicitly invoked, $a,b,c,x$ are
positive real numbers, $N,m$ are nonnegative integers, and $K$ is a positive
integer.  We use $H_0=0$ and
\[
 H_n^{(r)}=\sum_{j=1}^n\frac1{j^r},\qquad H_n=H_n^{(1)},
 \qquad (y)_N=\prod_{j=0}^{N-1}(y+j).
\]
Empty products equal one.  The falling factorial is
$\fall nN=n(n-1)\cdots(n-N+1)$.
The generalized Stieltjes constants have the convention
\begin{equation}\label{eq:stieltjes}
 \zeta(1+t,c)=\frac1t+
 \sum_{h=0}^\infty\frac{(-1)^h\gamma_h(c)}{h!}t^h,
 \qquad \gamma_0(c)=-\psi(c).
\end{equation}
We reserve $\zeta^{[q]}(s,c)=\partial_s^q\zeta(s,c)$ for derivatives in
\emph{order}; $\psi^{(q)}$ means a derivative in \emph{argument}.
A jet means an ordinary Taylor coefficient:
$[t^m]f(t)=f^{(m)}(0)/m!$.
This normalization matters in every all-order formula.

Our complete exponential Bell polynomials are specified by
\begin{equation}\label{eq:bell}
 \exp\left(\sum_{k\ge1}X_k\frac{t^k}{k!}\right)
 =\sum_{m\ge0}\Bcal_m(X_1,\ldots,X_m)\frac{t^m}{m!}.
\end{equation}
Thus $\Bcal_1=X_1$, $\Bcal_2=X_1^2+X_2$, and
$\Bcal_3=X_1^3+3X_1X_2+X_3$.
Reciprocal gamma is always interpreted as the entire function $1/\Gamma$,
not as a separately evaluated division by an infinite number.
